\providecommand{\Res}[2]{\qty[round-mode=figures,round-precision=3]{\fpeval{#1}}{#2}}
\providecommand{\Val}[1]{\num[round-mode=figures,round-precision=4]{\fpeval{#1}}}
\providecommand{\ResSci}[2]{\qty[round-mode=figures,round-precision=3,exponent-mode=scientific]{\fpeval{#1}}{#2}}
\subsubsection*{Question a) Angle critique}
\begin{enumerate}[label=\arabic*)]
    \item \textbf{Variables connues et recherchées :} $n_1 = \num{\v{n1}}$ (cœur), $n_2 = \num[minimum-decimal-digits=2]{\fpeval{\v{n1} - \v{dn}}}$ (gaine) ; variable recherchée : $\theta_c$
    \item \textbf{Équations pertinentes :}
    \[ n_1\sin\theta_c = n_2\sin(\qty{90}{\degree}) \]
    \item \textbf{Isolation de la variable recherchée :} 
    \[ \theta_c = \arcsin\left(\frac{n_2}{n_1}\right) \]
    \item \textbf{Résolution numérique :}
    \[ \theta_c = \arcsin\left(\frac{\num[minimum-decimal-digits=2]{\fpeval{\v{n1} - \v{dn}}}}{\num{\v{n1}}}\right) = \Res{asind((\v{n1} - \v{dn})/\v{n1})}{\degree} \]
    \item \textbf{Réponse :} $\theta_c = \mathbf{\Res{asind((\v{n1} - \v{dn})/\v{n1})}{\degree}}$. La réflexion totale n'est possible que si la lumière va vers un milieu moins réfringent : la gaine doit donc avoir un indice plus petit que le cœur.
\end{enumerate}
\subsubsection*{Question b) Vitesse dans le cœur}
\begin{enumerate}[label=\arabic*)]
    \item \textbf{Variables connues et recherchées :} $c = \qty{3.00e8}{m/s}$, $n_1 = \num{\v{n1}}$ ; variable recherchée : $v$
    \item \textbf{Équations pertinentes :}
    \[ n = \frac{c}{v} \]
    \item \textbf{Isolation de la variable recherchée :} 
    \[ v = \frac{c}{n_1} \]
    \item \textbf{Résolution numérique :}
    \[ v = \frac{\qty{3.00e8}{m/s}}{\num{\v{n1}}} = \ResSci{3e8/\v{n1}}{m/s} \]
    \item \textbf{Réponse :} $v = \mathbf{\ResSci{3e8/\v{n1}}{m/s}}$.
\end{enumerate}
\subsubsection*{Question c) Durée du trajet}
\begin{enumerate}[label=\arabic*)]
    \item \textbf{Variables connues et recherchées :} $L = \qty{\v{L}}{km} = \qty{\vsi{L}}{m}$, $v = \ResSci{3e8/\v{n1}}{m/s}$ ; variable recherchée : $t$ (trajet le plus court, en ligne droite)
    \item \textbf{Équations pertinentes :}
    \[ v = \frac{L}{t} \]
    \item \textbf{Isolation de la variable recherchée :} 
    \[ t = \frac{L}{v} \]
    \item \textbf{Résolution numérique :}
    \[ t = \frac{\qty{\vsi{L}}{m}}{\num[round-mode=figures,round-precision=4,exponent-mode=scientific]{\fpeval{3e8/\v{n1}}}\,\unit{m/s}} = \ResSci{\vsi{L}*\v{n1}/3e8}{s} = \Res{1e6*\vsi{L}*\v{n1}/3e8}{\micro\second} \]
    \item \textbf{Réponse :} l'impulsion met au moins \textbf{\Res{1e6*\vsi{L}*\v{n1}/3e8}{\micro\second}} (un peu plus en réalité, car elle zigzague).
\end{enumerate}
